% Figura: reparametrización de LoRA. Adaptada de hu2021lora (Figura 1).
% Se usa en sec:marco_llm_lora, después de la Ecuación de LoRA.

\begin{figure}[H]
    \centering
    \begin{tikzpicture}[x=1mm, y=1mm]

        \node[figentrada, text width=10mm, minimum height=7mm] (x) at (0, 0) {$\mathbf{x}$};

        % Rama congelada.
        \node[fignodo, text width=36mm, minimum height=24mm, fill=figrelleno] (W) at (-24, 25)
            {pesos preentrenados\\$W \in \mathbb{R}^{d\times k}$\\[2pt]{\footnotesize congelados}};

        % Rama entrenable: A reduce a r, B vuelve a d.
        \draw[draw=fignaranja!70!black, fill=fignaranja!25, thick] (8, 16) -- (32, 16) -- (26, 22) -- (14, 22) -- cycle;
        \node[font=\footnotesize] at (20, 19) {$A \in \mathbb{R}^{r\times k}$};
        \draw[draw=fignaranja!70!black, fill=fignaranja!25, thick] (14, 28) -- (26, 28) -- (32, 34) -- (8, 34) -- cycle;
        \node[font=\footnotesize] at (20, 31) {$B \in \mathbb{R}^{d\times r}$};
        \node[font=\footnotesize, text=fignaranja!70!black, align=left, anchor=west] at (34, 25)
            {entrenables\\rango $r \ll \min(d,k)$};

        % Suma y salida.
        \node[circle, draw=figtrazo, inner sep=1pt, font=\small] (sum) at (0, 46) {$+$};
        \node[figsalida, text width=26mm, minimum height=7mm] (h) at (0, 58) {$\mathbf{h} = W\mathbf{x} + BA\mathbf{x}$};

        \draw[figflecha] (x.north) -- ++(0, 4) -| (W.south);
        \draw[figflecha] (x.north) -- ++(0, 4) -| (20, 16);
        \draw[figflecha] (20, 22) -- (20, 28);
        \draw[figflecha] (W.north) |- (sum.west);
        \draw[figflecha] (20, 34) |- (sum.east);
        \draw[figflecha] (sum.north) -- (h.south);

    \end{tikzpicture}
    \caption{Reparametrización de LoRA con el incremento de bajo rango $BA$. Adaptada de \textcite{hu2021lora}.}
    \label{fig:marco_lora}
\end{figure}
